\chapter{ТЕОРЕТИЧЕСКИЕ ОСНОВЫ РАЗРАБОТКИ}

В первой главе рассматриваются предметная область, выбранная нормативная база и
математическая модель вычисления термодинамических свойств воды.
Основное внимание уделено структуре стандарта IAPWS-IF97 и способу расширения
расчётного маршрута по параметрам плотности и температуры.

\section{Анализ предметной области и существующих решений}

В инженерной практике свойства воды и водяного пара используются при
проектировании теплообменного оборудования, расчёте циклов паросиловых
установок, анализе режимов турбин и котельных агрегатов, а также при решении
задач автоматизации технологических процессов.
Во всех перечисленных случаях вычисления должны быть одновременно быстрыми,
повторяемыми и согласованными между разными программными интерфейсами.
Поэтому в качестве основы модели целесообразно использовать промышленный
стандарт, специально ориентированный на эффективные расчёты в широком диапазоне
давлений и температур \cite{iapws_if97_2012}.

Разрабатываемый программный комплекс должен обрабатывать несколько типовых пар
входных параметров: давление и температуру, давление и энтальпию, давление и
энтропию, давление и степень сухости, а также плотность и температуру.
Результат расчёта должен представлять не одну скалярную величину, а целостное
состояние рабочей среды, включающее давление, температуру, удельный объём,
плотность, энтальпию, энтропию, внутреннюю энергию, изобарную теплоёмкость,
скорость звука и номер региона IF97.

\begin{table}[htbp]
\centering
\caption{Основные входные и выходные параметры вычислительной модели}
\begin{tblr}{
colspec={|c|l|X|},
row{1}={font=\bfseries, halign=c},
hlines,
vlines
}
Обозначение & Параметр & Использование в модели \\
$p$ & Давление & Базовый независимый параметр для прямых и обратных маршрутов расчёта \\
$T$ & Температура & Используется для определения региона и вычисления производных термодинамического потенциала \\
$h$ & Энтальпия & Применяется в обратных задачах и в интерфейсах инженерного ввода \\
$s$ & Энтропия & Используется при расчётах по термодинамическим диаграммам и в обратных зависимостях \\
$\rho$ & Плотность & Необходима для маршрута $\rho T$ и анализа состояния в области высоких плотностей \\
$x$ & Степень сухости & Характеризует положение в двухфазной области на линии насыщения \\
\end{tblr}
\end{table}

Стандарт IAPWS-IF97 выбран потому, что он специально предназначен для
промышленных и расчётных приложений, где важна вычислительная эффективность.
Область допустимых значений разбита на несколько регионов, для каждого из
которых задан собственный набор уравнений.
Такое разбиение усложняет маршрутизацию вычислений, но позволяет достигать
высокой скорости работы и устойчивости результатов на границах применимости
модели \cite{iapws_if97_2012}.

\subsection{Выбор стандарта IAPWS-IF97 и границ области расчёта}

В рамках IF97 рабочая область разбивается на регионы 1--5 и область насыщения,
описываемую отдельным уравнением региона 4.
Регион 1 соответствует сжатой воде, регион 2 --- перегретому пару,
регион 3 --- высокоплотной околокритической области, регион 4 задаёт линию
насыщения, а регион 5 используется для высокотемпературного диапазона.
Дополнительно стандарт содержит формулу для метастабильного пара при давлениях
до 10 МПа \cite{iapws_if97_2012}.

Такое представление удобно для промышленного применения, но требует явной
классификации точки состояния перед вычислением итоговых свойств.
Следовательно, математическая модель дипломного проекта должна включать не
только сами уравнения состояния, но и топологию области расчёта: определение
региона, обработку границы насыщения и переходов через границу B23.

\section{Математическая модель расчёта термодинамических свойств}

В основе модели лежит представление термодинамического состояния через
безразмерные потенциалы, заданные различными уравнениями в зависимости от
региона.
Для регионов 1, 2 и 5 используется безразмерная энергия Гиббса
\cite{iapws_if97_2012}

\begin{equation}
\gamma(\pi,\tau) = \frac{g(p,T)}{RT},
\end{equation}

где $R$ --- газовая постоянная для воды и водяного пара,
$\pi = \frac{p}{p^{*}}$ --- безразмерное давление,
$\tau = \frac{T^{*}}{T}$ --- безразмерная температура.

Через производные функции $\gamma$ вычисляются основные свойства состояния:

\begin{equation}
v = \frac{RT}{p} \pi \gamma_{\pi}, \qquad
h = RT \tau \gamma_{\tau}, \qquad
s = R(\tau \gamma_{\tau} - \gamma),
\end{equation}

\begin{equation}
u = RT(\tau \gamma_{\tau} - \pi \gamma_{\pi}), \qquad
c_{p} = -R \tau^{2} \gamma_{\tau\tau}.
\end{equation}

Для региона 3, где более естественными независимыми переменными являются
плотность и температура, используется безразмерная энергия Гельмгольца
\cite{iapws_if97_2012,iapws_vpt_region3}

\begin{equation}
\phi(\delta,\tau) = \frac{f(\rho,T)}{RT}, \qquad
\delta = \frac{\rho}{\rho^{*}},
\end{equation}

и давление определяется выражением

\begin{equation}
p = \rho R T \delta \phi_{\delta}.
\end{equation}

Таким образом, стандарт обеспечивает единый физический базис для получения
плотности, удельного объёма, энтальпии, энтропии и других величин, но выбор
конкретного уравнения зависит от положения точки в фазовой области.

\subsection{Организация прямых и обратных маршрутов расчёта}

Прямой маршрут расчёта по паре $(p, T)$ начинается с валидации диапазонов
входных значений, после чего выполняется определение региона IF97.
На следующем шаге выбирается соответствующее уравнение состояния, и по его
производным восстанавливается полный набор выходных характеристик.
Для обратных маршрутов $(p, h)$, $(p, s)$ и $(p, x)$ применяются стандартные
обратные зависимости и уравнения насыщения, предусмотренные семейством
документов IAPWS.

С практической точки зрения важно, что программная модель не должна дублировать
формулы в каждом интерфейсе.
Все вычислительные маршруты необходимо свести к единому ядру, которое принимает
типизированный запрос и возвращает унифицированное описание состояния воды.
Такой подход упрощает тестирование, снижает риск расхождения результатов и
позволяет повторно использовать алгоритм в настольном приложении, web-оболочке
и gRPC-сервисе.

\subsection{Расширение модели методом секущих для маршрута rhoT}

Стандарт IF97 не задаёт универсального прямого расчёта по паре
независимых параметров $(\rho, T)$ во всей рабочей области в виде единого
закрытого выражения.
Поэтому в дипломном проекте используется численное расширение на основе метода
секущих.
Идея состоит в сведении задачи к поиску такого давления $p$, при котором
модель IF97 воспроизводит заданную плотность при фиксированной температуре.

Целевая функция записывается в виде

\begin{equation}
F(p) = \rho_{\mathrm{if97}}(p, T) - \rho_{\mathrm{target}}.
\end{equation}

Итерационный шаг метода секущих имеет вид

\begin{equation}
p_{k+1} =
p_{k} - F(p_{k}) \frac{p_{k} - p_{k-1}}{F(p_{k}) - F(p_{k-1})}.
\end{equation}

Итерации продолжаются до выполнения критерия по невязке или по изменению
давления между соседними шагами.
После нахождения давления дальнейший расчёт выполняется стандартным маршрутом
IF97 для пары $(p, T)$, благодаря чему весь набор итоговых свойств получается в
том же формате, что и для остальных входных режимов.

Применение метода секущих позволяет расширить практическую область использования
модели без нарушения её нормативной основы.
Стандартные региональные уравнения продолжают использоваться как первичный
источник термодинамических свойств, а численная часть отвечает только за
восстановление недостающего параметра в тех режимах, где это необходимо.
